% fuzz -t docs/display_restart_readiness.tex
% Reproduce all checks, including fidelity: bash scripts/verify_readiness_model.sh
\documentclass[a4paper,10pt,fleqn]{article}
\usepackage[margin=1in]{geometry}
\usepackage{fuzz}
\begin{document}
\title{Display Restart Readiness}
\author{A bounded model of the readiness observer}
\date{November 2026}
\maketitle

\section{Boundary and refinement}
The observer waits for a new process, then requires consecutive observations of
one Hub dial for the stability window. A process can keep its PID while the Hub
redials it. The generation is the pinned request port's opaque name, not a later
read of the registry. This models \texttt{DisplayReadiness},
\texttt{DisplayStability}, \texttt{ReadinessProbe}, and \texttt{ReadinessHTTP}.

The carrier has two non-reused dial names, a three-tick stability window, and a
six-tick total deadline. One tick conservatively permits more probe work than
the implementation's half-second cap. The smaller total deadline keeps exhaustive
exploration finite; it is not a claim to verify arbitrary real-valued clocks.
All time, including waits and individual blocking reads, is explicit. Zero-time
operations are local bookkeeping. The extra seventh tick is intentional: an
unbounded probe must be able to violate the deadline, not be disabled by a type.

\texttt{Snapshot} is the Hub reply's observation point. A later local ownership
check cannot credit its elapsed time to that reply. A disconnect after the last
sample cannot be ruled out by any observer; safety is about the age of the dial
at the accepted sample, not a promise that the process will remain alive later.
Kernel ownership and PID checks are abstracted as successful stage transitions
or \texttt{Crash}/\texttt{ProbeError}; real socket tests audit that refinement.
Failure-only logging and launchd diagnostics occur after this observer stops;
their existing separate bounded budget is not readiness work.

\section{Types and switches}
\begin{zed}
  ZBOOL ::= ztrue | zfalse \\
  RPHASE ::= rscandidate | rsstabilizing | rsready | rsfailed \\
  RWORK ::= rwwait | rwidle | rwlocal | rwhub | rwafter \\
  RPID ::= rpold | rpnew | rpgone \\
  TICK == 0 \upto 7 \\
  DIAL == 0 \upto 2
\end{zed}
\begin{axdef}
  RESETDIAL, BOUNDWORK : ZBOOL
\where
  RESETDIAL = ztrue \\
  BOUNDWORK = ztrue
\end{axdef}
Zero is no dial. Serial names increase even across \texttt{Drop}, excluding ABA
reuse. The script generates fidelity variants by changing exactly one switch.

\section{State}
As in \texttt{board\_ordering.tex}, safety is checked by reachability goals,
not inserted into the state predicate: doing the latter would disable the very
bad transitions the fidelity model must exhibit.
\begin{schema}{Readiness}
  phase : RPHASE \\
  work : RWORK \\
  process : RPID \\
  now, born, sampledBorn, sampledAt, started, age : TICK \\
  generation, serial, seen, remembered : DIAL
\end{schema}
\begin{schema}{Init}
  Readiness'
\where
  phase' = rscandidate \land work' = rwwait \land process' = rpold \\
  now' = 0 \land born' = 0 \land sampledBorn' = 0 \\
  sampledAt' = 0 \land started' = 0 \land age' = 0 \\
  generation' = 0 \land serial' = 0 \land seen' = 0
    \land remembered' = 0
\end{schema}

\section{Environment and polling}
\begin{schema}{Spawn}
  \Delta Readiness
\where
  phase \in \{rscandidate, rsstabilizing\} \land process = rpold \\
  process' = rpnew \land phase' = phase \land work' = work \\
  now' = now \land born' = born \land sampledBorn' = sampledBorn \\
  sampledAt' = sampledAt \land started' = started \land age' = age \\
  generation' = generation \land serial' = serial \\
  seen' = seen \land remembered' = remembered
\end{schema}
\begin{schema}{Dial}
  \Delta Readiness
\where
  phase \in \{rscandidate, rsstabilizing\} \land process = rpnew \\
  serial < 2 \land now < 6 \\
  serial' = serial + 1 \land generation' = serial' \land born' = now \\
  phase' = phase \land work' = work \land process' = process \\
  now' = now \land sampledBorn' = sampledBorn \land sampledAt' = sampledAt \\
  started' = started \land age' = age \\
  seen' = seen \land remembered' = remembered
\end{schema}
\begin{schema}{Drop}
  \Delta Readiness
\where
  phase \in \{rscandidate, rsstabilizing\} \land generation > 0 \\
  generation' = 0 \land serial' = serial \land born' = born \\
  phase' = phase \land work' = work \land process' = process \\
  now' = now \land sampledBorn' = sampledBorn \land sampledAt' = sampledAt \\
  started' = started \land age' = age \\
  seen' = seen \land remembered' = remembered
\end{schema}
\begin{schema}{Wait}
  \Delta Readiness
\where
  phase \in \{rscandidate, rsstabilizing\} \land work = rwwait \land now < 6 \\
  now' = now + 1 \land work' = rwidle \\
  (now' < 6 \land phase' = phase) \lor
    (now' = 6 \land phase' = rsfailed) \\
  process' = process \land born' = born \land sampledBorn' = sampledBorn \\
  sampledAt' = sampledAt \land started' = started \land age' = age \\
  generation' = generation \land serial' = serial \\
  seen' = seen \land remembered' = remembered
\end{schema}
\begin{schema}{NoCandidate}
  \Delta Readiness
\where
  phase = rscandidate \land work = rwidle \land process = rpold \\
  work' = rwwait \land phase' = phase \land process' = process \\
  now' = now \land born' = born \land sampledBorn' = sampledBorn \\
  sampledAt' = sampledAt \land started' = started \land age' = age \\
  generation' = generation \land serial' = serial \\
  seen' = seen \land remembered' = remembered
\end{schema}

\section{Probe stages}
\begin{schema}{Begin}
  \Delta Readiness
\where
  phase \in \{rscandidate, rsstabilizing\} \land work = rwidle \\
  process = rpnew \land now < 6 \land work' = rwlocal \\
  phase' = phase \land process' = process \land now' = now \\
  born' = born \land sampledBorn' = sampledBorn \land sampledAt' = sampledAt \\
  started' = started \land age' = age \\
  generation' = generation \land serial' = serial \\
  seen' = seen \land remembered' = remembered
\end{schema}
\begin{schema}{LocalDone}
  \Delta Readiness
\where
  phase \in \{rscandidate, rsstabilizing\} \land work = rwlocal \\
  work' = rwhub \land phase' = phase \land process' = process \land now' = now \\
  born' = born \land sampledBorn' = sampledBorn \land sampledAt' = sampledAt \\
  started' = started \land age' = age \\
  generation' = generation \land serial' = serial \\
  seen' = seen \land remembered' = remembered
\end{schema}
\begin{schema}{Snapshot}
  \Delta Readiness
\where
  phase \in \{rscandidate, rsstabilizing\} \land work = rwhub \\
  work' = rwafter \land seen' = generation \\
  sampledBorn' = born \land sampledAt' = now \\
  phase' = phase \land process' = process \land now' = now \land born' = born \\
  started' = started \land age' = age \\
  generation' = generation \land serial' = serial \land remembered' = remembered
\end{schema}
\begin{schema}{Finish}
  \Delta Readiness
\where
  phase \in \{rscandidate, rsstabilizing\} \land work = rwafter \land now < 6 \\
  process' = process \land now' = now \land born' = born \\
  sampledBorn' = sampledBorn \land sampledAt' = sampledAt \\
  generation' = generation \land serial' = serial \land seen' = seen \\
  (seen = 0 \land phase' = rscandidate \land remembered' = 0
    \land started' = 0 \land age' = 0) \\
  \quad~ \lor (seen > 0 \land remembered' = seen
    \land age' = sampledAt - sampledBorn \\
  \quad~ \land (((phase = rscandidate \lor
    (RESETDIAL = ztrue \land remembered \neq seen))
    \land started' = sampledAt) \\
  \quad~ \lor (phase = rsstabilizing \land
    (RESETDIAL = zfalse \lor remembered = seen) \land started' = started)) \\
  \quad~ \land ((sampledAt - started' \geq 3 \land phase' = rsready) \\
  \quad~ \lor (sampledAt - started' < 3 \land phase' = rsstabilizing))) \\
  (phase' = rsready \land work' = rwidle) \lor
    (phase' \neq rsready \land work' = rwwait)
\end{schema}
\begin{schema}{Spend}
  \Delta Readiness \\
  cost? : 1 \upto 2
\where
  phase \in \{rscandidate, rsstabilizing\} \\
  work \in \{rwlocal, rwhub, rwafter\} \land now < 6 \\
  (BOUNDWORK = ztrue \land cost? \leq 1 \land now + cost? \leq 6)
    \lor (BOUNDWORK = zfalse \land now + cost? \leq 7) \\
  now' = now + cost? \\
  (now' < 6 \land phase' = phase \land work' = work) \lor
    (now' \geq 6 \land phase' = rsfailed \land work' = rwidle) \\
  process' = process \land born' = born \land sampledBorn' = sampledBorn \\
  sampledAt' = sampledAt \land started' = started \land age' = age \\
  generation' = generation \land serial' = serial \\
  seen' = seen \land remembered' = remembered
\end{schema}
\texttt{Spend} can occur repeatedly in any probe stage. Thus a slow trickle of
HTTP header/body bytes is not collapsed into a single unbounded receive.

\section{Failure and completion}
\begin{schema}{Crash}
  \Delta Readiness
\where
  phase \in \{rscandidate, rsstabilizing\} \land process = rpnew \\
  phase' = rsfailed \land work' = rwidle \land process' = rpgone \\
  generation' = 0 \land serial' = serial \land now' = now \land born' = born \\
  sampledBorn' = sampledBorn \land sampledAt' = sampledAt \\
  started' = started \land age' = age \land seen' = seen
    \land remembered' = remembered
\end{schema}
\begin{schema}{ProbeError}
  \Delta Readiness
\where
  phase \in \{rscandidate, rsstabilizing\} \\
  work \in \{rwlocal, rwhub, rwafter\} \\
  phase' = rsfailed \land work' = rwidle \land process' = process \\
  now' = now \land born' = born \land sampledBorn' = sampledBorn \\
  sampledAt' = sampledAt \land started' = started \land age' = age \\
  generation' = generation \land serial' = serial \\
  seen' = seen \land remembered' = remembered
\end{schema}
\begin{schema}{Finished}
  \Xi Readiness
\where
  phase \in \{rsready, rsfailed\}
\end{schema}
The terminal self-loop is explicit: a finished observer has no further work.
Every nonterminal idle state enables \texttt{Begin} or \texttt{NoCandidate},
every wait enables \texttt{Wait}, and every probe stage enables its completion.
No synthetic error transition is needed to escape a blocked live state.

\section{Verification and partitions}
The verification script checks four safety invariants as unreachable goals:
ready implies sampled dial age at least three; time never exceeds six;
ready implies the new PID; ready occurs strictly before six. A separate model
check enables deadlock detection. Terminal self-loops represent termination, not
continued I/O. Removing \texttt{RESETDIAL} must reach a too-young ready dial;
removing \texttt{BOUNDWORK} must reach tick seven. See the companion coverage
audit for operation-to-code refinement and every test partition.
\end{document}
